% Part III of the merged paper: what could not be rerun, and why.
% Draft of 2026-10-06. Inputs after part2-published.tex.
% The R6 subsection has a slot for reading B, which is still collecting.

\section{What a rerun needs}
\label{sec:materials}

The seven targets were chosen because each released prompts, or data and code.
On four of them we could not assemble everything a rerun needs. This section
lists what was missing in each case, because it is the cost of Part~II and
because the same list is what we now ship with our own work.

The claim here is deliberately small. Reproducibility has been measured at
scale, on samples of hundreds and thousands of papers, and this is not that. Our
seven targets were selected for a different purpose, and seven cases support no
rate. What they can offer is a concrete list of the artifacts a rerun of this kind
turns out to depend on, which is more useful than a percentage.

\subsection{Five artifacts a rerun of this kind depends on}

\paragraph{Classifier weights.} The LABE metric is computed with three trained classifiers~\cite{wan2025}. The released folder carries configuration, scheduler and trainer state for each.
The weight files are not part of the public release, so we retrained all three
with their own script, reaching validation accuracies of 0.944, 0.933 and 0.941
against the reported 0.933, 0.933 and 0.917. Step~0 then came in at a median absolute relative difference of 0.286 against a
registered tolerance of 0.25, and that difference cannot be separated from
retraining noise: three classifiers of near-identical accuracy produce variances
on the same generations differing by up to a factor of four. The reported ordering, with the intersectional variance largest, appears in 14
of 18 per-classifier rows. That ordering may also be expected from an eight-cell
variance containing both the gender and the race spread, which is one reason the
neutral arm was designed for this target. With the original weights the
question the neutral arm was built for could be settled; without them it stays
open.

\paragraph{The input to a preprocessing tool.} Wilson and Caliskan code resumes and job descriptions to occupations with a
public coding tool~\cite{wilson2024}. A rerun needs to know which text was sent
to that tool, and we could not determine it. We coded 4{,}859 distinct texts
through that tool's API and tried seven readings of the resume side, three
registered before the first attempt and four after it failed and before they
were computed. None of the seven reproduces the published counts, and the difference has the
same shape every time, concentrated in three of the nine occupations. The
consistency of that shape is what suggests a single preprocessing choice we did
not recover, rather than scattered noise. The job-description side reproduces under the first reading, 8 of 9 counts
within 20 percent, so what we were missing is specifically the resume-side
input.

\paragraph{A runtime version.} Ghanbari Haez and Dragoni generated their answers through a local model
server~\cite{ghanbari2025}. Which version of it produced the released answers
is something we were unable to establish, and generation is sensitive to it. Our
regeneration missed the registered fidelity tolerance by 0.0006. A serial rerun,
registered as a diagnostic before it was run, did not close the gap: their step
is 0.8583, our parallel run 0.8680, our serial run 0.8721. Request batching is
therefore not the explanation, and no conclusion about their claim is drawn. What our own generation does show is
reported separately and is a statement about our answers only.

\paragraph{How the comparison groups were formed.} ABLEIST reports compounding
harm across stacked attributes~\cite{ableist2025}. Its Table~6 caption leaves
several groupings open, and all three of our registered readings came in at best
6 of 12 numbers per model against a rule of 10. One mixed reading, which we
reconstructed by search over the readings the caption leaves open, reproduces
all 24 numbers exactly on both models: the three single-and-double columns are
unrestricted, covering every gender and both castes, while the triple-stacked
column is restricted to the marginalised values. Since that reading reproduces every published number to the second decimal on
two independent models, it is very likely the computation behind the table,
though we cannot establish that from the released materials alone. It also
indicates that the composition of the comparison groups may influence the
observed increase, and that this is worth considering when interpreting that
row. The reconstruction did not satisfy the preregistered reproduction criterion, and
the reconstructed interpretation is reported alongside that outcome.

\paragraph{The frame a fine-tuned model expects.} The same work releases a
fine-tuned labeller. Running it needs the instruction string it was trained
with, which we inferred from the released inference code. The labeller could not
in the end be run at the scale our check needed, which is a compute limit on our
side and is reported as one.

\subsection{Beck et al.: how close the reproduction comes}

Beck et al. report that combining all sociodemographic attributes produces the
largest change in predictions for most of their datasets~\cite{beck2024}. Their
code, their sampling notebook and six small open models make this the one target
that is fully runnable, and we ran their pipeline unchanged over six models and
seven conditions.

The registered reading required the printed winner to be ours and our combined
share to be within five percentage points, on at least four of six models.
Three of the six meet both conditions. On a fourth the combination is also the
most influential condition, so the direction of the claim holds on four of the
six, but our share for that model is 13.02 against 7.00 and so misses the
five-point band. The two remaining models put a single attribute first instead:
one changes between 0.0 and 0.1 percent of labels in every condition, so its
ordering turns on a single label, and the other is inside the five-point band on
size while ordering gender first by 2.1 points.

The preregistration names a fallback reading, the same comparison averaged over
all three of their prompt formats, to be used if the first fails. \textbf{It is not reported here, because its collection was not finished in
time.} The first
format is complete and is what the paragraph above rests on; the other two are
partial. No number from them appears in this paper, and the fallback reading
stays open rather than being decided on part of its data.

What the first format does carry is a check on this work rather than on theirs.
The collection was lost to a deleted working directory in October and rebuilt
from nothing: a new environment at a new path from the versions the authors pin,
the source data downloaded again, the sample rebuilt from their notebook. All
thirty-six cells of the recollection reproduce the first to the second decimal,
and the verdict is unchanged. The distilled predictions are in the repository
with a manifest of checksums, so the table above can be recomputed without the
raw collection.

\subsection{What this section is for}

It is not a reproducibility audit and it is not a judgement on any of these
studies, none of which is shown here to be wrong. Of the five artifact types listed above, two could be reconstructed or inferred
from what was released, with different levels of confidence and neither fully
verified: the grouping behind the published cells, which reproduces every
published number but was found by search after the registered readings did not
match and which did not meet the preregistered criterion, and the frame a
fine-tuned model expects, which we inferred from the released inference code but
could not verify at the scale our check needed. The other three could not, which
is worth knowing in both directions. The point is forward-looking and applies first to us: a
rerun of an interaction claim turns out to depend on a short and predictable
list, namely the trained weights behind any learned metric, the exact input to
any preprocessing tool, the runtime version used for generation, the grouping
behind each published cell, and the frame a fine-tuned model expects. Part~I and Part~II of this paper ship each of these that applies to them, and
that is the habit we would encourage.
